\section{Conclusion}
\label{sec:conclusion}

This survey has examined the rapidly evolving landscape of AI coding agents, from inline code completion to fully autonomous software development systems.

\subsection{Summary of Contributions}

We have presented:
\begin{itemize}
    \item A taxonomy categorizing coding agents by autonomy level, interaction scope, and integration mode
    \item Systematic coverage of major systems across four categories: code completion tools, conversational assistants, autonomous agents, and integrated development environments
    \item Analysis of evaluation benchmarks including HumanEval, MBPP, and SWE-bench, with discussion of their limitations
    \item Identification of open challenges in reliability, security, code quality, and human-agent collaboration
\end{itemize}

\subsection{Key Takeaways}

Several themes emerge from our analysis:

\textbf{Rapid capability growth}: The field has progressed from code completion to autonomous agents in approximately two years, with systems now capable of resolving real-world software issues.

\textbf{Persistent reliability challenges}: Despite impressive demonstrations, agents remain unreliable on complex, long-horizon tasks. Error rates that seem acceptable on benchmarks compound in practice.

\textbf{Evaluation gaps}: Current benchmarks, while valuable, do not fully capture real-world utility. Code quality, security, and maintainability remain difficult to measure.

\textbf{Research opportunities}: The intersection of AI and software engineering offers rich opportunities for research in human-AI collaboration, evaluation methodology, and specialized agent development.

\subsection{Looking Forward}

Coding agents represent a fundamental shift in software development. The tools surveyed here are early steps in a trajectory toward more capable AI assistance. The software engineering research community has a critical role to play---not only in advancing these capabilities, but in ensuring they improve rather than degrade software quality, developer skills, and the craft of programming.

We hope this survey serves as a useful map of the current landscape and a foundation for future research.
